\chapter{Marco Teorico y Referencial}\label{chapter:marco_teorico_y_referencial}

\section{Estado del Arte}

Históricamente, para la generación de informes médicos el uso de lápiz y papel fue la norma.
Aunque la digitalización inicial buscaba reducir la carga administrativa, 
la transcripción manual de datos se convirtió en un nuevo obstáculo para los profesionales de la salud,
limitando la agilidad de los sistemas modernos.

\subsection{\englishText{Automatic Speech Recognition} (\englishText{ASR})}
Con la llegada de los modelos de reconocimiento automático de voz (\englishText{ASR}) se ha abierto un nuevo 
camino para la generación de informes médicos. Estos modelos son capaces de transcribir
el habla en texto, pero estos presentaron fallas importantes:

\boldparagraph{Falta de especialización}{
    dificultad para entender terminología médica técnica.
}

\boldparagraph{Alta tasa de error}{
    los sistemas básicos llegaban a tener altos porcentajes de errores.
}

\boldparagraph{Carga de revisión}{
    debido a estos errores, los médicos debían dedicar mucho tiempo a revisiones manuales para garantizar 
    la seguridad del paciente.
}

\subsection{Automatización con \englishText{ASR} y Modelos de Lenguaje (\englishText{LLM})}

El estado actual de la técnica se enfoca en la integracion de \englishText{ASR} y Modelos de Lenguaje (\englishText{LLM}),
logrando automatizar de forma \englishText{end-to-end} la generación de informes médicos como se 
hace mencion en estos trabajos recientes:   
\citeText{adedeji2024soundhealthcareimprovingmedical, li2024searchingbestpracticesmedical} ambos
haciendo uso del modelo de \englishText{Whisper} de \englishText{OpenAI} y de diversas herramientas y tecnicas para mejorar la transcripcion de los audios, 
logrando asi reducir la tasa de error de palabra (\englishText{WER}) y la carga de revisiones manuales.
Por otro lado,\citeText{10912421} propone utilizar realizar \englishText{fine-tuning} del modelo 
\englishText{Whisper} para el mejoramiento del de la tasa de error de palabra (\englishText{WER}) y la utilizacion 
de servicos de terceros para convertir los conceptos médicos extraídos a vocabularios 
médicos estándar (\englishText{CIE-10-ES} y \englishText{SNOMED CT})

A pesar de los avances presentados por \citeText{adedeji2024soundhealthcareimprovingmedical, li2024searchingbestpracticesmedical} y \citeText{10912421}, 
se observa una limitación: la dependencia casi exclusiva de la capacidad generativa intrínseca del \englishText{LLM} para estructurar el informe clínico
final a partir de la transcripción cruda. Si bien la reducción del \englishText{WER} mediante técnicas de \englishText{fine-tuning} mejora la fidelidad
del texto, la síntesis del informe sigue estando sujeta a la interpretación del modelo, lo que incrementa el riesgo de omisiones o alucinaciones en
datos clínicos críticos.

Para mitigar esta vulnerabilidad, se propone introducir una capa intermedia basada en Procesamiento
del Lenguaje Natural (\englishText{PLN}) especializado que implementa el Reconocimiento de Entidades Nombradas (\englishText{NER})
para identificar de forma explícita los conceptos médicos presentes en el discurso. 

Una vez extraídas, estas entidades no solo se mapean a vocabularios estándar de interoperabilidad (\englishText{CIE-10-ES} y \englishText{SNOMED CT}), 
sino que se utilizan de forma activa como llaves de consulta para extraer información clínica estructurada adicional.
Al alimentar al \englishText{LLM} no solo con el audio transcrito, sino con un contexto enriquecido de metadatos clínicos normalizados y verificados,
se reduce la incertidumbre del modelo generativo. Esto permite mayor nivel de detalle de informes clínicos generados,
precisión diagnóstica y rigurosidad técnica.

\section{Marco Teorico}

% Dentro de este trabajo fin de master se tomaran en cuenta los siguientes conceptos teoricos:

\subsection{Aplicacion web}

Una aplicación web es un sistema de software (sistema distribuido del tipo cliente-servidor) basado
en tecnologías y estándares del Consorcio \englishText{World Wide Web} (\englishText{W3C}, por sus siglas en inglés) 
que proporciona recursos web específicos, como contenido y servicios, a través de la interfaz
de usuario provista por el navegador web.

\subsection{Aplicacion Movil}

Una aplicación móvil es un software diseñado para ejecutarse en dispositivos móviles, 
tales como smartphones y tabletas. Estas aplicaciones se desarrollan para aprovechar las capacidades específicas de los dispositivos móviles, 
como la cámara, el \englishText{GPS}, el acelerómetro y la pantalla táctil.

Dentro del contexto de este trabajo fin de master, la aplicación móvil desemepñará el rol de cliente grabando 
el audio del dictado médico.

\subsection{Aplicación \englishText{Backend}}
El \englishText{backend} representa la capa de lógica de negocio, procesamiento de datos 
e integración de servicios en una arquitectura cliente-servidor. A diferencia de la aplicación móvil
(\englishText{frontend}), el \englishText{backend} opera en el entorno del servidor y se encarga de 
gestionar la persistencia de datos, la autenticación, la autorización y la orquestación de los servicios,
entre otros. 
En el contexto de este sistema, actúa como el núcleo central que recibe el audio de la 
aplicación móvil, coordina los pipelines de procesamiento lingüístico y estructuración, 
y crea los informes clínicos en formato estructurado.


\subsection{\englishText{SSL/TLS}}
Los protocolos \englishText{Secure Sockets Layer} (\englishText{SSL}) y su sucesor evolutivo, 
\englishText{Transport Layer Security} (\englishText{TLS}), son protocolos criptográficos diseñados para garantizar
la seguridad, confidencialidad e integridad de la información transmitida a través de una red informática. 

\englishText{TLS} opera en la capa de transporte del modelo \englishText{TCP/IP}, situándose entre la aplicación 
(\englishText{HTTP} o \englishText{WebSocket}) y el protocolo de red. De acuerdo con \citeText{rfc8446}, su funcionamiento se basa en dos
fases principales: el emparejamiento criptográfico inicial (\englishText{handshake}), donde el cliente y el servidor autentican sus 
identidades mediante certificados digitales y negocian los algoritmos de cifrado, y la posterior fase de 
transferencia de datos protegida por dicha negociación.

\subsection{\englishText{WebSocket}}
\englishText{WebSocket} es un protocolo de comunicación bidireccional sobre la capa de transporte
\englishText{TCP} que permite el intercambio de datos en tiempo real entre un cliente y un servidor
\citeText{rfc6455}. 
A diferencia del modelo tradicional de petición-respuesta de \englishText{HTTP}, el protocolo establece un apretón de manos
(\englishText{handshake}) inicial para actualizar la conexión, permitiendo que ambas partes envíen mensajes
de forma continua de datos básico y sin la sobrecarga de renegociar la sesión en cada intercambio.

En el contexto de este sistema, la aplicación móvil inicia la comunicación con el servidor \englishText{backend}
mediante un \textit{handshake} \englishText{HTTP}. Si la negociación es exitosa, el protocolo se actualiza a \englishText{WebSocket} seguro
(\englishText{WSS}), estableciendo un canal cifrado mediante \englishText{SSL/TLS} sobre el cual circulan los mensajes en tiempo real 
necesarios para el flujo de audio.

\begin{figure}[htbp]
    \centering
    \begin{tikzpicture}[
        font=\sffamily,
        box/.style={draw, thick, minimum width=4.5cm, minimum height=7cm, align=center},
        arrow/.style={-{Stealth[scale=1.2]}, thick},
        biarrow/.style={{Stealth[scale=1.2]}-{Stealth[scale=1.2]}, thick},
        lbl/.style={align=center, font=\small, text width=5.8cm} 
    ]

    \node[box] (mobileBox) at (0,0) {};
    
    % Título Superior
    \node[align=center, font=\large\bfseries, anchor=north] at (0,3.2) {Dispositivo Móvil\\(\englishText{Smartphone})};
    
    % Icono de Smartphone
    \begin{scope}[shift={(0,0)}]
        \draw[thick, rounded corners=6pt] (-0.5,-0.8) rectangle (0.5,0.8);
        \draw[thick, rounded corners=2pt] (-0.4,-0.6) rectangle (0.4,0.6);
        % Botón de inicio / Línea inferior
        \draw[thick] (-0.15,-0.7) -- (0.15,-0.7);
        % Icono interior (Estilo App Store)
        \draw[thick] (-0.2,-0.2) -- (0,-0.5) -- (0.2,-0.2);
        \draw[thick] (-0.25,-0.25) -- (0.25,-0.25);
        \node[font=\tiny\bfseries] at (0,0.1) {A};
    \end{scope}

    \node[box] (serverBox) at (11,0) {};
    
    % Título Superior
    \node[align=center, font=\large\bfseries, anchor=north] at (11,3.2) {Servidor \englishText{Backend}};
    
    % Icono de Servidor (Rack) - Ajustado el shift para la nueva posición
    \begin{scope}[shift={(10.4,-0.4)}]
        \foreach \y in {0, 0.3, 0.6, 0.9} {
            \draw[thick, rounded corners=1pt, fill=white] (0,\y) rectangle (1.2,\y+0.2);
            \fill (0.15,\y+0.1) circle (1.2pt);
            \fill (0.3,\y+0.1) circle (1.2pt);
            \draw (0.45,\y+0.1) -- (1.05,\y+0.1);
        }
        % Icono de Nube superpuesto
        \node[cloud, draw, thick, fill=white, cloud puffs=8, cloud puff arc=120, minimum width=1.1cm, minimum height=0.7cm] at (1.1,0.2) {};
    \end{scope}
    
    
    % Coordenadas de los bordes internos recalculadas para X=11
    \coordinate (L1) at (2.25, 2.3);
    \coordinate (R1) at (8.75, 2.3);
    
    \coordinate (L2) at (2.25, 0.8);
    \coordinate (R2) at (8.75, 0.8);
    
    \coordinate (L3) at (2.25, -0.7);
    \coordinate (R3) at (8.75, -0.7);
    
    \coordinate (L4) at (2.25, -2.2);
    \coordinate (R4) at (8.75, -2.2);

    % Flecha 1: Petición Inicial
    \draw[arrow] (L1) -- (R1) node[lbl, midway, above] {Petición de Conexión WSS (Inicial)};

    % Flecha 2: Intercambio de Datos
    \draw[biarrow] (L2) -- (R2) 
        node[lbl, midway, above] {Intercambio de Datos}
        node[lbl, midway, below] {Bidireccional (\englishText{WSS})};

    % Flecha 3: Handshake (Derecha a Izquierda)
    \draw[arrow] (R3) -- (L3) 
        node[lbl, midway, above] {Mensaje de Confirmación}
        node[lbl, midway, below] {(\englishText{Handshake})};

    % Flecha 4: Canal Cifrado
    \draw[biarrow] (L4) -- (R4) 
        node[lbl, midway, above] {Canal \englishText{WSS} Cifrado (\englishText{SSL/TLS})}
        node[lbl, midway, below] {Transmisión Segura};

    \end{tikzpicture}
    \caption{Diagrama de arquitectura y flujo de conexión \englishText{WSS}.}
    \label{fig:diagrama_wss}
\end{figure}


\subsection{\englishText{Gateway}}
El \englishText{gateway} o puerta de enlace es un dispositivo o software que actúa como interfaz de conexión entre 
sistemas o redes que utilizan distintos protocolos de comunicación. En nuestra implementación, esta pasarela
funciona como intermediario entre la aplicación móvil y el \englishText{backend} a través del protocolo
\englishText{WSS}. Su función principal es recibir los audios enviados por la aplicación móvil, 
transformarlos y transmitirlos al \englishText{backend} para su posterior procesamiento.


\subsection{Reconocimiento Automático de Voz (\englishText{ASR})}
% El Reconocimiento Automático de Voz (\englishText{Automatic Speech Recognition} o \englishText{ASR})
% es una disciplina de la inteligencia artificial cuyo objetivo es la conversión de una señal acústica
% de habla humana en un texto equivalente legible por una máquina. 
% De forma matemática, el proceso se modela como la búsqueda de la secuencia de palabras 
% $\hat{W}$ que maximiza la probabilidad condicionada dada una secuencia de observaciones acústicas $X$ \citeText{jm3}:

% $$\hat{W} = \arg\max_{W} P(W|X)$$

% Los sistemas modernos de \englishText{ASR} han migrado de arquitecturas tradicionales basadas en 
% Modelos Ocultos de Márikov (\englishText{HMM}) hacia redes neuronales profundas del tipo \englishText{end-to-end}
% (E2E), como las arquitecturas basadas en \englishText{Transformers}. 
% Estos sistemas procesan el audio mediante representaciones espectrográficas (frecuencia y tiempo) y 
% utilizan mecanismos de autoatención (\englishText{self-attention}) para capturar de manera simultánea 
% las dependencias acústicas y las relaciones semánticas del lenguaje natural, permitiendo procesar el 
% contexto global del dictado clínico.

El Reconocimiento Automático de Voz (\englishText{Automatic Speech Recognition} o \englishText{ASR}) es 
una disciplina de la inteligencia artificial cuyo objetivo es la conversión de una señal acústica de habla
humana en un texto equivalente legible.

Desde una perspectiva funcional, un sistema de \englishText{ASR} moderno actúa como un modelo de caja negra que opera en tres etapas fundamentales \citeText{jm3}:
\begin{itemize}
    \item \textbf{Entrada:} El sistema recibe una señal de audio (el habla humana), la cual se transforma internamente en representaciones basadas en frecuencia y tiempo.
    \item \textbf{Procesamiento:} El modelo analiza de forma simultánea tanto las características acústicas del sonido como las relaciones semánticas del lenguaje, evaluando probabilísticamente qué secuencias de palabras tienen mayor sentido dentro del contexto.
    \item \textbf{Salida:} Produce como resultado final una transcripción y estructurada del dictado.
\end{itemize}

Para nuestra implementación, se utilizarán como parte fundamental del procesamiento de audios para poder convertirlo
en texto plano.

\begin{figure}[htbp]
    \centering
    \begin{tikzpicture}[
        font=\sffamily,
        box/.style={draw, thick, align=center, inner xsep=20pt, inner ysep=15pt},
        arrow/.style={-{Stealth[scale=1.2]}, thick},
        neuron/.style={circle, draw, fill=white, minimum size=4pt, inner sep=0pt},
        lbl/.style={align=center, font=\small}
    ]

    % Generación limpia de la señal de Audio analógica (Posición fija a la izquierda)
    \begin{scope}[shift={(0,0)}]
        \foreach \x/\h in {
            -1.6/0.3, -1.5/0.7, -1.4/0.5, -1.3/1.2, -1.2/0.8, -1.1/1.5, -1.0/0.9, -0.9/2.0, 
            -0.8/1.4, -0.7/0.6, -0.6/1.1, -0.5/0.8, -0.4/1.6, -0.3/0.9, -0.2/1.3, -0.1/0.7, 
             0.0/1.7,  0.1/1.1,  0.2/1.4,  0.3/0.6,  0.4/1.2,  0.5/0.5,  0.6/1.0,  0.7/1.8,
             0.8/1.3,  0.9/0.7,  1.0/1.5,  1.1/0.9,  1.2/1.1,  1.3/0.4,  1.4/1.2,  1.5/0.6,
             1.6/0.8
        } {
            \draw[thick, blue!60!black] (\x, -\h/2) -- (\x, \h/2);
        }
        \draw[thick, gray] (-1.7, 0) -- (1.7, 0);
    \end{scope}

    % Caja central (ASR) con su propio texto interno directamente en el nodo para que el padding funcione
    \node[box] (asrBox) at (5,0) {
        Motor de Reconocimiento\\
        Automático de Voz (\englishText{ASR})
    };
    
    % Títulos Superiores
    \node[font=\large\bfseries, align=center] (inputTitle) at (0, 2.5) {Entrada de Audio};
    \node[font=\large\bfseries, align=center] (outputTitle) at (10, 2.5) {Salida de Texto};
    
    % Representación de texto transcrito (Derecha)
    \node[align=left, font=\small, text width=3.5cm] (outputTxt) at (10.2, 0) {
        \texttt{``El paciente presenta}\\[2pt]
        \texttt{dolor torácico ...``}\\[2pt]
    };

    % Flecha Izquierda: Va desde el final de la onda (1.7,0) hasta el borde oeste de la caja padded
    \draw[arrow] (1.7, 0) -- (asrBox.west) node[lbl, midway, above=3pt] {};
    
    % Flecha Derecha: Va desde el borde este de la caja padded hasta el inicio del texto de salida
    \draw[arrow] (asrBox.east) -- (8.4, 0) node[lbl, midway, above=3pt] {};

    \end{tikzpicture}
    \caption{Diagrama funcional de un sistema de Reconocimiento Automático de Voz (\englishText{ASR}).}
    \label{fig:diagrama_asr}
\end{figure}

\subsection{Procesamiento del Lenguaje Natural (\englishText{PLN})}
El Procesamiento del Lenguaje Natural (\englishText{Natural Language Processing} o \englishText{PLN}) 
es una ramade la ciencia de la computación, la lingüística y la inteligencia artificial orientada 
al desarrollo de metodologías algorítmicas que permiten a los ordenadores comprender, 
interpretar y manipular el lenguaje humano de forma automatizada \citeText{jm3}.

El \englishText{PLN} abarca tareas complejas que van desde el análisis morfosintáctico elemental 
(como la tokenización y el etiquetado gramatical) hasta la comprensión del lenguaje natural 
(\englishText{NLU}), la gestión del conocimiento semántico y el reconocimiento de entidades nombradas (\englishText{NER}). 
Este último permite extraer de forma automática identificar entidades en categorias previamente definidas.

\subsection{Reconocimiento de Entidades Nombradas (\englishText{NER})}
El Reconocimiento de Entidades Nombradas (\englishText{Named Entity Recognition} o \englishText{NER}) 
constituye una subtarea fundamental del \englishText{PLN} cuyo propósito es localizar y clasificar 
elementos de texto no estructurado en categorías semánticas previamente definidas \citeText{jm3}.

Los modelos \englishText{NER} se pueden entrenar para especializarse en un dominio específico, que en nuestro caso es 
el dominio médico, se entrenan para escanear las transcripciones médicas e identificar y categorizar 
entidades médicas de interés como patologías, procedimientos quirúrgicos, fármacos, anatomía y variables fisiológicas. 
Las arquitecturas contemporáneas de \englishText{NER} emplean modelos basados en representaciones de 
codificadores lingüísticos profundos (como variantes de \englishText{BERT} o \englishText{RoBERTa}) 
adaptadas al sector biomédico, los cuales analizan el contexto izquierdo y derecho de cada palabra
para clasificar los límites de una entidad médica con alta precisión.

\begin{figure}[htbp]
    \centering
    \begin{tikzpicture}[
        font=\sffamily,
        % Caja central con padding generoso usando inner sep
        box/.style={
            draw, 
            thick, 
            align=center, 
            inner xsep=20pt, 
            inner ysep=30pt,
            minimum height=4.5cm
        },
        arrow/.style={-{Stealth[scale=1.2]}, thick},
        lbl/.style={align=center, font=\large\bfseries}
    ]

    % =========================================================================
    % 1. BLOQUE CENTRAL (Posición fija en el origen)
    % =========================================================================
    \node[box] (plnBox) at (0,0) {
        \Large\bfseries Procesamiento de Lenguaje\\
        \Large\bfseries Natural (PLN) y NER
    };

    % =========================================================================
    % 2. BLOQUE IZQUIERDO (Texto de Entrada) - Desplazado a la izquierda x=-5.5
    % =========================================================================
    \node[align=center, text width=4.5cm] (inputGroup) at (-5.5, 0) {
        % Icono de documento dibujado nativamente
        \begin{tikzpicture}
            \draw[thick, rounded corners=2pt] (-0.4, -0.5) rectangle (0.4, 0.6);
            \draw[thick] (-0.2, 0.3)  -- (0.2, 0.3);
            \draw[thick] (-0.2, 0.1)  -- (0.2, 0.1);
            \draw[thick] (-0.2, -0.1) -- (0.2, -0.1);
            \draw[thick] (-0.2, -0.3) -- (0.0, -0.3);
        \end{tikzpicture} \\[10pt]
        % Frase de entrada
        \large Hoy, María visitó la tienda de Apple en Madrid para comprar un iPhone.
    };
    
    % Título del bloque izquierdo (Situado arriba del grupo)
    \node[lbl] (inputTitle) at (-5.5, 2.8) {Texto de Entrada};


    % =========================================================================
    % 3. BLOQUE DERECHO (Resultado de NER) - Desplazado a la derecha x=5.8
    % =========================================================================
    \node[align=left, text width=6.5cm] (outputTxt) at (5.8, 0) {
        \large [María]\textbf{\color{blue!80!black}{Persona}} visitó la tienda 
        de [Apple]\textbf{\color{orange!80!black}{Organización}} en 
        [Madrid]\textbf{\color{green!50!black}{Ubicación}} para 
        comprar un [iPhone]\textbf{\color{red!70!black}{Producto}}.
    };
    
    % Título del bloque derecho (Alineado perfectamente con el izquierdo)
    \node[lbl] (outputTitle) at (5.8, 2.8) {Resultado de NER};


    % =========================================================================
    % 4. CONEXIONES ENLAZADAS A LOS BORDES DE LA CAJA (AUTOMÁTICAS)
    % =========================================================================
    % Flecha desde el extremo derecho del grupo de entrada hasta el borde izquierdo de la caja
    \draw[arrow] (inputGroup.east) -- (plnBox.west);
    
    % Flecha desde el borde derecho de la caja hasta el extremo izquierdo del texto procesado
    \draw[arrow] (plnBox.east) -- (outputTxt.west);

    \end{tikzpicture}
    \caption{Diagrama funcional corregido para el proceso de NER.}
    \label{fig:diagrama_ner_corregido}
\end{figure}

\subsection{Modelo de Lenguaje de Gran Escala (\englishText{LLM})}

Un Modelo de Lenguaje de Gran Escala (\englishText{Large Language Model} o \englishText{LLM}) es un sistema 
de inteligencia artificial basado en redes neuronales profundas que ha sido entrenado mediante aprendizaje 
auto-supervisado sobre corpus de texto masivos. Su arquitectura subyacente se fundamenta en el mecanismo 
de \englishText{Transformer} \citeText{vaswani2017attention}, lo que les permite capturar dependencias lingüísticas
de largo alcance y demostrar capacidades emergentes complejas como el razonamiento en contexto, la 
traducción de estilos narrativos y la comprensión de instrucciones estructuradas \citeText{wei2022emergent}.

En el contexto de esta aplicación, el \englishText{LLM} no se utiliza como una fuente primaria de diagnóstico, 
sino como el redactor de la información extraída. Su propósito fundamental es actuar como un transformador
de formato: toma la transcripción literal del audio médico, las entidades médicas extraídas 
junto a los detalles extraidos de los vocabularios estándar de interoperabilidad (\englishText{SNOMED CT} y 
\englishText{CIE-10-ES}) para generar un informe clínico y coherente.



\subsection{Vocabularios Estándar de Interoperabilidad (\englishText{CIE-10-ES} y \englishText{SNOMED CT})}
La interoperabilidad semántica en los sistemas de información sanitaria digital exige el empleo de terminologías
de referencia y clasificaciones normalizadas a nivel internacional. En este sistema se consideran dos estándares
fundamentales:

\boldparagraph{SNOMED CT}{
    La \englishText{Systematized Nomenclature of Medicine - Clinical Terms} es la terminología clínica integral
    y polijerárquica más semánticamente rica del mundo. Organizada mediante conceptos, descripciones y relaciones
    lógicas, permite codificar con granularidad cualquier juicio clínico, hallazgo, estructura anatómica
    o procedimiento. Facilita que diferentes sistemas informáticos interpreten con exactitud inequívoca el 
    significado de los términos registrados por el personal clínico.
}

\boldparagraph{CIE-10-ES}{
    Es la Clasificación Internacional de Enfermedades, 10.ª revisión, correspondiente a la Modificación
    Clínica adoptada de manera oficial para la codificación diagnóstica y de procedimientos en el Sistema
    Nacional de Salud de España. A diferencia de \englishText{SNOMED CT} (que es una terminología analítica), 
    la \englishText{CIE-10-ES} es un sistema de clasificación cerrado y estadístico, esencial para la gestión
    epidemiológica, la facturación sanitaria, la indexación de historias clínicas y la auditoría médica
    institucional.
}

Ambos forman parte del sistema que se quiere desarrollar desempeñando el rol de proveedores y validadores
de información clínica. Donde al identificar una entidad médica, se puede: buscar esta entidad, validar su existencia
y obtener su información detallada y obtener un identificador único para su uso de caracter nacional e internacional,
\englishText{SNOMED CT} y \englishText{CIE-10-ES}.


\section{Marco Conceptual}

El diseño conceptual de este sistema se fundamenta en la transición de un flujo de información puramente secuencial y probabilístico hacia un modelo híbrido de conocimiento guiado. El objetivo principal es mitigar el riesgo inherente de alucinación semántica de los modelos generativos mediante la validación y el enriquecimiento ontológico previo.

\subsection{Flujo Arquitectónico de Datos y Conocimiento}
Como se ilustra en la Figura \ref{fig:diagrama_conceptual}, el ciclo de vida del dato clínico sigue un pipeline lineal y determinista en sus fases tempranas, transformándose en un proceso de enriquecimiento contextual antes de la intervención del modelo generativo:

\begin{enumerate}
    \item \textbf{Captura e Ingesta (Audio):} La voz del profesional de la salud es transmitida en tiempo real desde la aplicación móvil hacia el backend mediante el protocolo \englishText{WebSocket}.
    \item \textbf{Transcodificación Acústica (ASR):} El modelo \englishText{Whisper} procesa la señal de audio cruda y genera una transcripción textual textual (\englishText{raw text}), la cual posee un alto riesgo de contener errores de sintaxis médica o ruido conversacional.
    \item \textbf{Extracción Semántica (PLN/NER):} El texto plano es procesado por un pipeline de PLN especializado, aislando exclusivamente las entidades clínicas de interés (síntomas, diagnósticos, fármacos) y descartando conectores lingüísticos irrelevantes.
    \item \textbf{Normalización y Envejecimiento Óntico (SNOMED CT / CIE-10-ES):} Las entidades extraídas son mapeadas a códigos únicos de los vocabularios estándar. Este paso actúa como un validador de veracidad clínica; si un concepto no se puede normalizar o validar contra la ontología, se marca para revisión, eliminando la ambigüedad lingüística.
    \item \textbf{Inyección de Contexto y Síntesis (LLM):} El \englishText{LLM} no recibe la transcripción cruda a ciegas. En su lugar, se le alimenta con un \englishText{prompt} estructurado que contiene: la transcripción original más los metadatos estructurados y verificados de las ontologías. El modelo actúa estrictamente como un formateador avanzado, generando el informe clínico final con alta fidelidad y nulo margen de alucinación.
\end{enumerate}

% NOTA: Recuerda definir este label en tu entorno de figura TikZ, Mermaid o Gráfico externo
% \label{fig:diagrama_conceptual}

\subsection{Definición de Informe Enriquecido}
Dentro del marco de esta investigación, se define formalmente un \textbf{Informe Enriquecido} como el artefacto clínico textual resultante que supera cualitativa y cuantitativamente a una transcripción médica estándar. Mientras que una transcripción tradicional es un reflejo literal (y potencialmente erróneo) del discurso hablado, el \textit{Informe Enriquecido} se caracteriza por poseer las siguientes propiedades operacionales:

\begin{itemize}
    \item \textbf{Estructuración Formal:} La información se organiza bajo estándares clínicos reconocidos (por ejemplo, el formato SOAP: Subjetivo, Objetivo, Análisis, Plan).
    \item \textbf{Metadatos Semánticos Incrustados:} Cada diagnóstico, síntoma o procedimiento clave del texto está indexado inequívocamente a sus correspondientes códigos \englishText{SNOMED CT} y \englishText{CIE-10-ES}.
    \item \textbf{Densidad Informativa Optimizada:} Se elimina el lenguaje coloquial, las vacilaciones del habla y las redundancias, manteniendo una alta rigurosidad técnica y precisión diagnóstica.
\end{itemize}